\documentclass{article}
\usepackage[utf8]{inputenc}
\usepackage[T1]{fontenc}
\usepackage[margin=2cm]{geometry}
\usepackage{hyperref}
\sloppy
\begin{document}
\section*{Historical source}
Encyclopaedia Britannica scan, source row 2171. Parse 3.5 extracted the text. Super 3.5 VL corrected it from an enlarged crop of each block. It can still contain recognition errors.
Page headers, footers, and page numbers are removed.

\subsection*{ATARAXY (E0185)}
ATARAXY, a term used by the stoics and sceptics, to denote that calmness of mind which secures us from all emotions arising from vanity and self-conceit.

ATARNEĀ, an ancient town of Myſia, ſituated between Adramyttium and Pitane, remarkable for the marriage of Ariſtotle with the ſiſter or concubine of the tyrant Hermias; alſo for the dotage of that philoſopher.

\subsection*{ATAXY (E0186)}
ATAXY, in a general sense, the want of order: With physicians, it signifies irregularity of crises and paroxysms of fevers.

\subsection*{ATCHE (E0187)}
ATCHE, in commerce, a small silver coin used in Turkey, and worth only one third of the English penny.

\subsection*{ATCHIEVEMENT (E0188)}
ATCHIEVEMENT, in heraldry, denotes the arms of a person, or family, together with all the exterior ornaments of the field; as helmet, mantle, crest, scrolls, and motto, together with such quarterings as may have been acquired by alliances, all marshalled in order.

\subsection*{ATCHIEVE (E0189)}
ATCHIEVE. This term is derived from the French achever, i. e. to finish or make an end of; but signifies, in heraldry, to perform great actions or exploits.

\subsection*{ATE (E0190)}
ATE, the goddess of mischief, in the Pagan theology. She was daughter of Jupiter, and cast down from heaven at the birth of Hercules. For Juno having deceived Jupiter, in causing Eurysthenes to be born before Hercules, Jupiter expelled his resentment on Ate, as the author of that mischief; and threw her headlong from heaven to earth, swearing she should never return thither again, (Homeri Il. six. 125.) The name of this goddess comes from ατα, noceo, to hurt. Her being the daughter of Jupiter means, according to mythologists, that no evil happens to us but by the permission of providence; and her banishment to earth denotes the terrible effects of divine justice among men.

\subsection*{ATEGUA (E0191)}
ATEGUA, or ATEGUA, an ancient town of Spain, placed by fame in the road from Antiquara, now Antequera, to Hifpalis, or Seville; by others near Alcala Real; which last is the more probable situation, because the flumen falum, now the Salado, was in its neighbourhood. Now Tebala Vieja, or Teivela.

\subsection*{ATELLA (E0192)}
ATELLA, an ancient town of Campania in Italy, between Capua and Neapolis. From this town the Atellanae fabulae, or Atellani ludi, took their name. These were also called Oseci, from their inventor, in whose territory Atella lay. They were generally a species of farce, interlarded with much ribaldry and buffoonery; and sometimes were exodia or interludes presented between the acts of other plays. The actors in these farces were not reckoned among the common players, nor deemed infamous; but retained the rights of their tribe, and might be listed for soldiers, the privilege only of free men. The ruins of this town are still to be seen about 11 miles from the modern Aversa, which was built out of its materials.

A TEMPO giusto, in music, signifies to sing or play in an equal, true, and just time.

\subsection*{ATERNUM (E0193)}
ATERNUM, a town of Lucania in Italy, now Aterni, (Cluvenus): Also a town in the territory of the Piceni, now Pescara, a port-town of Naples, situated on the Adriatic. E. Long. 15. 25. N. Lat. 42. 30.

ated on the Adriatic. E. Long. 15. 25. N. Lat. 42. 30.
ATESTE, a town in the territory of Venice in
Italy, now called Este. E. Long. 12. 6. N. Lat. 45. 25.

\subsection*{ATHAMADULET (E0194)}
ATHAMADULET, the prime minister of the Persian empire, as the grand vizier is of the Turkish empire. He is great chancellor of the kingdom, president of the council, superintendent of the finances, and is charged with all foreign affairs.

\subsection*{ATHAMANTA, SPIGNEL (E0195)}
ATHAMANTA, SPIGNEL; a genus of the di-nynia order, belonging to the pentandria class of plants. Of this genus Linnaeus enumerates nine species; but none of them merit particular notice except the cretensis, otherwise called daucus creticus. This is an umbelliferous plant growing wild in the Levant and the warmer parts of Europe. The leaves are irregularly difpoled, and formed like those of fennel. The flower-stalk rises about two feet high, sending out many branches, terminated at the top by compound umbels, composed of near 20 small ones. These have white flowers with five petals, which are succeeded by oblong, hairy, channelled fruit, divided into two parts, containing one oblong hairy feed. The feeds have a warm biting taste, with an agreeable aromatic smell. They are kept in the shops as a medicine, are carminative, and said to be diuretic; but are little used in practice. The plant may be propagated from feeds, which should be sown on an open bed of light dry ground; the following autumn the plants should be taken up, and planted at about a foot distance in a bed of light sandy earth, where the roots will continue several years.

\subsection*{ATHANASIA, GOLDLOCKS (E0196)}
ATHANASIA, GOLDLOCKS; a genus of the polygamia æqualis order, belonging to the syngenesia class of plants. Of this genus, Dr. Linnaeus mentions 11

Species; of which the following fix are by Mr Miller reckoned worthy of a place in those gardens where there are conveniences for preferring tender plants, as none of the species will outlive the winters of this country. 1. The dentata grows naturally at the Cape of Good Hope. It rises three or four feet high, sending out many side-branches, garnished with pretty long narrow leaves, having some resemblance to those of buckthorn plantain. The branches are generally terminated by yellow flowers formed into a kind of umbel. These flowers appear early in summer, and the seeds ripen in autumn. 2. The trifurcata is also a native of Africa, as are the three following sorts. It is a shrubby plant, about the height of the former, sending out several weak shoots, garnished with awl-shaped trifid leaves set close to them, of a pale silvery colour. The flowers are yellow, and terminate the branches in a corymbus. They appear in August, and the seeds ripen in October. 3. The crithmifolia rises with a strong shrubby stalk to the height of six or eight feet, sending out many ligneous branches garnished with long narrow leaves terminating in four or five parts like those of samphire. The flowers are like those of the former, and appear in July and August, being succeeded by seeds which ripen in October. 4. The pubescens hath strong woody stalks four or five feet high, covered with a woolly bark, as are also the branches, which are closely garnished with entire, spear-shaped, woolly leaves, and are terminated by yellow flowers having a long foot-stalk, and are sometimes succeeded by seeds that ripen in the autumn. 5. The annua is a low annual plant, seldom rising above a foot high, and sending forth two or three slender branches garnished with wing-pointed leaves: the stalks are terminated


\section*{Source record}
\begin{description}
\item[Scan ID] \texttt{britannica:309f765b1ecc83e0a7f7762e44b4a927a7555443:train:2171}
\item[Image file] \texttt{page-2171-f8869fa95f19.jpg}
\item[Image SHA-256] \texttt{f8869fa95f190d0e556bf8d2397c0c36da436ea1b23405995f0aac0f398f26d5}
\item[Dataset] NationalLibraryOfScotland/encyclopaedia\_britannica\_illustrated at revision \texttt{309f765b1ecc83e0a7f7762e44b4a927a7555443} (CC0-1.0)
\item[Extraction] model \texttt{nvidia/nemotron-3.5-parse}, raw reply \texttt{extract/raw/page-2171-f8869fa95f19.json}
\item[Raw reply SHA-256] \texttt{bdad01469c86af1aa0c878fb933f1193633bedfcea0f8956eb6fa31a05c04ab1}
\item[Correction] model \texttt{nvidia/Nemotron-3.5-Super-VL}, block crops, raw replies \texttt{correct/raw/crop\_nothink/page-2171-f8869fa95f19-b*.json}; super\_vl: 16
\item[Running header] A S T [ 852] A S T
\end{description}
\end{document}
